\documentclass[12pt]{article}
\usepackage[T1]{fontenc}
\usepackage[utf8]{inputenc}
\usepackage[brazil]{babel}
\usepackage[margin=1in]{geometry}
\usepackage{newtxtext,newtxmath}
\usepackage{ragged2e}
\usepackage[normalem]{ulem}
\setlength{\parindent}{0pt}
\linespread{1.15}
\setlength{\parskip}{0.7\baselineskip}
\emergencystretch=3em
\sloppy
\begin{document}
\justifying
\noindent \textbf{Processo n.: }\uline{5017817-46.2025.4.02.5101}

\noindent Verifica-se que a questão objeto da controvérsia jurídica suscitada no recurso manejado pela parte recorrente esteve afetada no representativo da controvérsia - \textbf{Tema n. }\textbf{16} da TNU - (PU no processo n. PEDILEF 0501999-48.2009.4.05.8500/ SE),  afetado em 12/09/2011 foi julgado (trânsito em julgado em 26/03/2015 - no STF), por meio do qual foi elucidada a seguinte questão:

{\leftskip=3.1cm\rightskip=0.8cm\parindent=0pt \textit{"Saber qual o termo inicial da progressão funcional de policial federal."}\par}

\noindent Restou firmada a seguinte tese:

{\leftskip=3.1cm\rightskip=0.8cm\parindent=0pt \textit{"A eficácia da progressão funcional deve ser observada segundo a situação individual de cada servidor e seus efeitos retroagem ao momento em que os requisitos legais foram implementados. Vide Tema 82 e PEDILEF n. 500367784.2014.404.7101. Entendimento revisado pelo PEDILEF n. 05207128420124058300 (orientação alinhada ao STJ: RESP 1.649.269/RJ)." (Tema 16 TNU).}\par}

\noindent Com efeito, depreende-se da leitura do voto condutor do acórdão que o mesmo se encontra em conformidade com o decidido no  tema pela TNU.

\noindent A proposito, consta do acordao hostilizado: (...) Gratificação de incentivo, Sistema Remuneratório e Benefícios, Servidor Público Civil, DIREITO ADMINISTRATIVO E OUTRAS MATÉRIAS DE DIREITO PÚBLICO

\noindent Nessas condições, o conteúdo do Acórdão recorrido é consentâneo com o entendimento da TNU, firmado em julgamento de recurso representativo de controvérsia, o que atrai a incidência da regra prevista pelo art. 14, III, \textit{b,} da Resolução CJF n. 586/2019:

\noindent \textit{\textbf{Art. 14.}}\textit{ Decorrido o prazo para contrarrazões, os autos serão conclusos ao magistrado responsável pelo exame preliminar de admissibilidade, que deverá, de forma sucessiva: (...) }\textit{\textbf{III }}\textit{- negar seguimento a pedido de uniformização de interpretação de lei federal interposto contra acórdão que esteja em conformidade com entendimento consolidado: }\textit{\textbf{b)}}\textit{ em }\textit{\textbf{\uline{recurso representativo de controvérsia pela Turma Nacional de Uniformização}}}\textit{ ou em pedido de uniformização de interpretação de lei dirigido ao Superior Tribunal de Justiça;(...).}

\noindent Posto isso, com arrimo no \textbf{\uline{art. 14, III, }}\textit{\textbf{\uline{b}}}\textit{,} da Resolução CJF n. 586/2019 c/c art. 1.030, I, do CPC, \textbf{NEGO SEGUIMENTO }\textbf{a}o\textbf{ PU}\textbf{.}

\noindent Intimem-se. Aguarde-se o decurso do prazo recursal. Após, certifique-se o trânsito em julgado e devolva-se ao juízo de origem.
\end{document}
